\documentclass{article}

\usepackage{microtype}
\usepackage{graphicx}
\usepackage{subcaption}
\usepackage{booktabs} % for professional tables
\usepackage{hyperref}
\newcommand{\theHalgorithm}{\arabic{algorithm}}
% Course style (ethdl2026, from Moodle); [accepted] shows the author names.
\usepackage[accepted]{ethdl2026}
\usepackage{amsmath}
\usepackage{amssymb}
\usepackage{mathtools}
\usepackage{amsthm}
\usepackage[capitalize,noabbrev]{cleveref}
\theoremstyle{plain}
\newtheorem{theorem}{Theorem}[section]
\newtheorem{proposition}[theorem]{Proposition}
\newtheorem{lemma}[theorem]{Lemma}
\newtheorem{corollary}[theorem]{Corollary}
\theoremstyle{definition}
\newtheorem{definition}[theorem]{Definition}
\newtheorem{assumption}[theorem]{Assumption}
\theoremstyle{remark}
\newtheorem{remark}[theorem]{Remark}
\usepackage[textsize=tiny]{todonotes}
\icmltitlerunning{ETHZ Deep Learning 2026 -- Project 14}

\begin{document}

\twocolumn[
  \icmltitle{ETHZ Deep Learning 2026 Project Proposal \\
    Project 14: How Much of the Learned Testing Policy Should a Model Keep?}

  \icmlsetsymbol{equal}{*}

  % TODO: full names and ETH usernames
  \begin{icmlauthorlist}
    \icmlauthor{Daniele Bugatti (\texttt{username})}{ETH}
    \icmlauthor{Rafael van 't Veer (\texttt{username})}{ETH}
    \icmlauthor{Mehdi Ejtehadi (\texttt{username})}{ETH}
    \icmlauthor{Wesley Remy (\texttt{weremy})}{ETH}
  \end{icmlauthorlist}

  \icmlaffiliation{ETH}{Department of Computer Science, ETH Zurich, Zurich, Switzerland}

  % The course style requires at least one corresponding author (a team member)
  \icmlcorrespondingauthor{Wesley Remy}{weremy@student.ethz.ch}

  \icmlkeywords{Observation shift, informative missingness, clinical time series, state-space models, domain adaptation}

  \vskip 0.3in
]

\printAffiliationsAndNotice{Advisor TA: Sarah Martinson. Topic: Project 14, Deep Time-Series Models under Observation Shift. }

\begin{abstract}
Clinical time-series models learn from \emph{when} and \emph{whether} tests are ordered, not only from their values. This testing pattern reflects one hospital's observation policy and changes when the model is deployed elsewhere. Before deployment, hospitals typically run a model silently for 1--3 months, which yields only a handful of labelled outcome cases. We propose a ladder of adaptations for a latent state-space model with a separate testing head: first re-learning that head from the new hospital's records, which needs no outcome labels, then the testing and outcome heads, then the full model. We ask which observation shifts can be corrected without labels, and how many outcome cases the others need.
\end{abstract}

\section{Introduction \& Related Work}

\textbf{Problem.} Deep models for clinical time series learn from what is measured and also from when it is measured: a lactate order signals that a doctor suspects sepsis. This testing pattern is predictive \citep{sisk2021}, but it reflects the hospital's observation policy, and it changes when a model trained at hospital A is deployed at hospital B (\emph{observation shift}). The change comes in two kinds. Hospital B may test differently given what the model can see: troponin is ordered for 3\% of patients at Beth Israel but 32\% at Emory (PhysioNet 2019, \citealp{reyna2020}), so an order means less there. Or B's clinicians may test on information the model never sees, so an order carries outcome information that B's records alone cannot reveal.

Before going live, models are often run silently for 1--3 months \citep{sendak2020}. That run-in yields the only labelled data from B, and very little of it: a typical ICU admits 600--1{,}300 patients a year \citep{checkley2014}, so a 1--3 month run-in contains about 2--6 sepsis cases. Outcomes are scarce, but B's records are not: they show which tests B orders for which patients. We therefore propose a \textbf{ladder of adaptations} for a latent state-space model in which the testing evidence enters through its own head: (1) re-learning the testing head from B's records without outcomes, the rest frozen; (2) re-learning the testing and outcome heads with the run-in's outcomes; (3) fine-tuning the whole model, starting from (1) and penalised towards it. Higher rungs repair more kinds of shift but need more outcome cases; rung 1 needs none.

\textbf{Research question.} \emph{Which observation shifts can be corrected without outcome labels by re-learning only the testing part, and how many outcome cases do the higher rungs need for the rest?} Concretely: (i) does rung 1 beat the fixed rules of prior work (always keep or always remove the pattern) and full fine-tuning at realistic run-in sizes; (ii) for which kinds of shift, and from how many outcome events, do rungs 2 and 3 pay off; (iii) does the state-space model adapt better than GRU-D, whose testing part cannot be separated, fine-tuned at its last layer or in full.

\textbf{Deep models use the observation pattern but do not separate it.} Architectures for irregular series consume masks and time gaps (GRU-D, \citealp{che2018}), continuous-time paths (Neural CDE, \citealp{kidger2020}; Latent ODE, \citealp{rubanova2019}) or time-value sets with attention (mTAN, \citealp{shukla2021}; Raindrop, \citealp{zhang2022}), and fuse the pattern with physiology in one representation. Models with a separate observation component exist, a Hawkes process for clinicians' measurement events \citep{alaa2017} and DeepJoint's observation heads \citep{jeanselme2022}, but none adapts that component alone at a new site.

\textbf{Observation shift between hospitals.} Multi-site deep sepsis and ICU benchmarks lose 0.08--0.09 AUROC across sites \citep{moor2023,yaib2024} and include missingness indicators without attributing the loss to them. Studies that isolate the pattern choose all or nothing: \citet{futoma2021} fitted with measurement indicators and dropped them at deployment, which transferred best from Beth Israel to a Memphis hospital but not in reverse; \citet{rockenschaub2024} found that ignoring informative missingness costs at the source but is more robust elsewhere, especially when missingness is highly informative. \citet{lasko2024} argue conceptually for a site-specific process model next to a stable disease model; \citet{zhou2023} adapt to missingness shift without labels in linear models. None re-learns the observation part of a deep sequence model at a new site.

\textbf{Adapting with few labels.} Clinical model updating is a ladder from recalibration to full refitting, and parsimonious updates win at small samples \citep{steyerberg2004,vergouwe2017}; a full refit is risky well below the 100 events needed just to validate a model \citep{riley2024}. In deep learning, which part to fine-tune depends on the kind of shift \citep{lee2023}, and retraining only the last layer can suffice \citep{kirichenko2023}, but clinical fine-tuning studies use hundreds to thousands of target patients \citep{moor2023,yaib2024}. \textbf{Gap:} no work re-learns only the observation part of a clinical sequence model from unlabelled target records, and none compares such adaptations across run-in sizes and kinds of shift.

\section{Experimental Setup \& Methods}

\textbf{Model.} A neural latent state-space model, extending our pilot: a discrete latent patient state ($K{\approx}32$ states) evolves by a learned transition, and three heads read it: a value head (likelihood of each measured value), a testing head (probability that each test is ordered, given the state) and an outcome head. Predictions come from filtering the state over time, trained end to end on the outcome loss plus the likelihood of values and tests. The testing pattern enters only through the testing head. Rung 1 re-learns the testing head by maximising the likelihood of B's values and test orders, with initial distribution, transition, value and outcome heads frozen; rung 2 also re-learns the outcome head with the run-in's outcomes; rung 3 fine-tunes everything. Baselines: keep and remove the testing evidence (weight $w{=}1$ and $w{=}0$, Futoma), a single tuned weight $w$ (our pilot), recalibration, and GRU-D, unchanged or fine-tuned at the last layer or in full.
\textbf{E1: simulation with known optimum.} Our latent-state simulator, whose exact Bayes predictor is known, with three shift types: B tests more or less selectively or at random (visible in B's records); B's clinicians also test on a hidden risk factor (visible only through outcomes); and the same patient state carries a different risk at B. We measure each rung's regret against the Bayes predictor at B over 1--104 weeks of deployment (up to about 50 events). Expected: rung 1 suffices for the first type; the others need rungs 2--3 and enough events.

\textbf{E2: real hospitals (dataset to be determined).} Requirements: several hospitals with hospital identifiers, time-stamped labs and vitals with the testing pattern preserved, and enough outcome events per site; the main candidate is eICU-CRD (208 US hospitals, \citealp{pollard2018}; credentialed access). Design: ordered source $\to$ target pairs, run-ins of 1, 3, 6, 12 and 24 months drawn from the target site (20 draws each), evaluation on all remaining target patients, primary metric AUROC (pre-specified), secondary calibration.

\textbf{E3: architecture.} Whether the state-space model adapts from fewer cases than GRU-D fine-tuned at the last layer or in full, on the pairs of E2.

\textbf{Computational cost and feasibility.} We have one NVIDIA GTX 1080 Ti (11\,GB) on snowflake and a 4-core laptop. All models are small (32 latent states or GRU hidden size 32--64, at most 48 hourly steps per patient), so GPU memory is not a constraint. Measured on the laptop CPU: the pilot state-space model trains in 15 minutes on 8{,}000 simulated patients, and 108 simulation fits of GRU, GRU-D, Neural CDE and Transformer took 3.6 hours in total; E1 therefore runs on the laptop. E2 and E3 need about 50 source models (sites $\times$ models $\times$ 3 seeds) and a few thousand short adaptation runs, which we run on the GPU. We will time one site first and, if the full design does not fit the available weeks, reduce the number of draws or pairs.

\section{Results \& Discussion}

\textbf{Preliminary results} (linear models on open PhysioNet 2019 and 2012 data; simulation). On Beth Israel $\to$ Emory, keeping the pattern was best, because there was barely any observation shift between the two (pattern-model transfer gap 0.03 AUROC, 95\% CI $-0.01$ to $0.07$); dropping it cost 0.05 AUROC. Across 26 unit pairs the pattern's meaning shifted clearly in 7, mostly for mortality; there, a single weight tuned on a 3--24 month run-in gained up to 0.02--0.03 AUROC and stayed within 0.001 elsewhere (after removing impossible lab values). A one-parameter update beat a full refit up to about 20 cases; the refit won from about 50. In simulation, rung 1 from 25 unlabelled patients came close to the exact optimum under four testing shifts (regret 0.0001--0.0003 Brier, keep up to 0.0012), and full fine-tuning was worse up to two years; with a hidden risk factor, only fine-tuning helped, after two years.

\section{Conclusion}

The project turns the keep-or-remove question of prior work into re-learning the testing part from the new hospital's own records, with outcome-based adaptation only where the run-in supports it. The simulator, the pilot model and the run-in pipeline already exist. Timeline: weeks 1--2 the multivariate state-space model (with static covariates) and E1; weeks 3--4 E2 and E3; weeks 5--6 analysis and report.

\bibliography{bibliography}
\bibliographystyle{ethdl2026}

%%%%%%%%%%%%%%%%%%%%%%%%%%%%%%%%%%%%%%%%%%%%%%%%%%%%%%%%%%%%%%%%%%%%%%%%%%%%%%%
%%%%%%%%%%%%%%%%%%%%%%%%%%%%%%%%%%%%%%%%%%%%%%%%%%%%%%%%%%%%%%%%%%%%%%%%%%%%%%%
% APPENDIX
%%%%%%%%%%%%%%%%%%%%%%%%%%%%%%%%%%%%%%%%%%%%%%%%%%%%%%%%%%%%%%%%%%%%%%%%%%%%%%%
%%%%%%%%%%%%%%%%%%%%%%%%%%%%%%%%%%%%%%%%%%%%%%%%%%%%%%%%%%%%%%%%%%%%%%%%%%%%%%%
% \newpage
% \appendix
% \onecolumn
% \section{Appendix}

\end{document}
